\documentclass[10pt,a4paper]{amsart}
\usepackage{amsmath,amssymb,amsthm,graphicx}
\usepackage{verbatim}
\newenvironment{reference}{\comment}{\endcomment}
%\newenvironment{reference}{}{}
\newenvironment{slogan}{\comment}{\endcomment}
\newenvironment{history}{\comment}{\endcomment}

% For commutative diagrams we use Xy-pic
\usepackage[all]{xy}

% We use 2cell for 2-commutative diagrams.
\xyoption{2cell}
% Theorem environments.
%
\theoremstyle{plain}
\newtheorem{theorem}[subsection]{Theorem}
\newtheorem{proposition}[subsection]{Proposition}
\newtheorem{lemma}[subsection]{Lemma}

\theoremstyle{definition}
\newtheorem{definition}[subsection]{Definition}
\newtheorem{example}[subsection]{Example}
\newtheorem{exercise}[subsection]{Exercise}
\newtheorem{situation}[subsection]{Situation}

\theoremstyle{remark}
\newtheorem{remark}[subsection]{Remark}
\newtheorem{remarks}[subsection]{Remarks}

\numberwithin{equation}{subsection}

% Macros
%
\def\lim{\mathop{\mathrm{lim}}\nolimits}
\def\colim{\mathop{\mathrm{colim}}\nolimits}
\def\Spec{\mathop{\mathrm{Spec}}}
\def\Hom{\mathop{\mathrm{Hom}}\nolimits}
\def\Ext{\mathop{\mathrm{Ext}}\nolimits}
\def\SheafHom{\mathop{\mathcal{H}\!\mathit{om}}\nolimits}
\def\SheafExt{\mathop{\mathcal{E}\!\mathit{xt}}\nolimits}
\def\Sch{\mathit{Sch}}
\def\Mor{\mathop{\mathrm{Mor}}\nolimits}
\def\Ob{\mathop{\mathrm{Ob}}\nolimits}
\def\Sh{\mathop{\mathit{Sh}}\nolimits}
\def\NL{\mathop{N\!L}\nolimits}
\def\CH{\mathop{\mathrm{CH}}\nolimits}
\def\proetale{{pro\text{-}\acute{e}tale}}
\def\etale{{\acute{e}tale}}
\def\QCoh{\mathit{QCoh}}
\def\Ker{\mathop{\mathrm{Ker}}}
\def\Im{\mathop{\mathrm{Im}}}
\def\Coker{\mathop{\mathrm{Coker}}}
\def\Coim{\mathop{\mathrm{Coim}}}

% Boxtimes
%
\DeclareMathSymbol{\boxtimes}{\mathbin}{AMSa}{"02}

%
% Macros for moduli stacks/spaces
%
\def\QCohstack{\mathcal{QC}\!\mathit{oh}}
\def\Cohstack{\mathcal{C}\!\mathit{oh}}
\def\Spacesstack{\mathcal{S}\!\mathit{paces}}
\def\Quotfunctor{\mathrm{Quot}}
\def\Hilbfunctor{\mathrm{Hilb}}
\def\Curvesstack{\mathcal{C}\!\mathit{urves}}
\def\Polarizedstack{\mathcal{P}\!\mathit{olarized}}
\def\Complexesstack{\mathcal{C}\!\mathit{omplexes}}
% \Pic is the operator that assigns to X its picard group, usage \Pic(X)
% \Picardstack_{X/B} denotes the Picard stack of X over B
% \Picardfunctor_{X/B} denotes the Picard functor of X over B
\def\Pic{\mathop{\mathrm{Pic}}\nolimits}
\def\Picardstack{\mathcal{P}\!\mathit{ic}}
\def\Picardfunctor{\mathrm{Pic}}
\def\Deformationcategory{\mathcal{D}\!\mathit{ef}}
\usepackage{cleveref}
\crefname{originallemma}{original lemma}{original lemmas}
\Crefname{originallemma}{Original lemma}{Original lemmas}
\makeatletter
\newlabel{algebra-lemma-hathat-finitely-generated}{{05GG}{source}{}{sourcepointer-1}{}}
\newlabel{algebra-lemma-hathat-finitely-generated@cref}{{[originallemma][0][]05GG}{[1][0][]source}{}{}{}}
\makeatother
\crefname{originallemma}{original lemma}{original lemmas}
\Crefname{originallemma}{Original lemma}{Original lemmas}
\makeatletter
\newlabel{algebra-lemma-quotient-complete}{{031A}{source}{}{sourcepointer-2}{}}
\newlabel{algebra-lemma-quotient-complete@cref}{{[originallemma][0][]031A}{[1][0][]source}{}{}{}}
\makeatother
\crefname{originallemma}{original lemma}{original lemmas}
\Crefname{originallemma}{Original lemma}{Original lemmas}
\makeatletter
\newlabel{formal-spaces-lemma-closed}{{0AMS}{source}{}{sourcepointer-3}{}}
\newlabel{formal-spaces-lemma-closed@cref}{{[originallemma][0][]0AMS}{[1][0][]source}{}{}{}}
\makeatother
\title{Closed submodules of I-adically complete modules: (1) closed iff (2) N = intersection (N + I\textasciicircum{}n M) iff (3) M/N complete; finite generation gives N complete}
\author{The Stacks Project Authors}
\date{stacks-project:examples.tex:local-snapshot-20261012; GNU FDL 1.2 or later}
\begin{document}
\maketitle
\noindent\textit{Editorial note.} All delivered mathematical text is the author's own, copied from the pinned Stacks Project chapter file examples.tex; no statement or proof step was rewritten, reordered or removed.\par\medskip
\noindent\textit{Editorial note.} The counted proof is the author's complete argument: The author marks where finite generation is used ('this is where we use that I is finitely generated') - preserved verbatim.\par\medskip
\noindent\textit{Editorial note.} The label lemma-closed-is-complete names the counted statement, and the proof environment immediately following it in the source is delivered verbatim as its proof.\par\medskip
\noindent\textit{Editorial note.} The counted statement is carried by the author's own lemma environment in the chapter source.\par\medskip
\noindent\textit{Allowed standard prerequisites (editorial).}\begin{itemize}
\item I-adic completions N\textasciicircum{} of submodules and density of N in N\textasciicircum{}; factorization of continuous maps into complete modules
\item Formal Spaces, closed lemma for (1) iff (2); Algebra, quotient-complete lemma for (2) iff (3); Algebra, hathat-finitely-generated lemma for direct summands of complete modules
\end{itemize}\medskip
\ifdefined\hypertarget\hypertarget{sourcepointer-1}{}\fi
\noindent\textit{Source reference.} Original lemma 05GG: sources/stacks-chow/tags.txt, source lines 7060--7060. Prerequisite pointer invoked at 836; it names a standard fact of a sibling chapter that the counted argument uses, and every step of the counted proof built on it is retained verbatim in the delivered body.\par\medskip
\ifdefined\hypertarget\hypertarget{sourcepointer-2}{}\fi
\noindent\textit{Source reference.} Original lemma 031A: sources/stacks-chow/tags.txt, source lines 3977--3977. Prerequisite pointer invoked at 824; it names a standard fact of a sibling chapter that the counted argument uses, and every step of the counted proof built on it is retained verbatim in the delivered body.\par\medskip
\ifdefined\hypertarget\hypertarget{sourcepointer-3}{}\fi
\noindent\textit{Source reference.} Original lemma 0AMS: sources/stacks-chow/tags.txt, source lines 13415--13415. Prerequisite pointer invoked at 822; it names a standard fact of a sibling chapter that the counted argument uses, and every step of the counted proof built on it is retained verbatim in the delivered body.\par\medskip
\begin{lemma}
\label{lemma-closed-is-complete}
Let $I \subset R$ be an ideal of a ring. Let $M$ be an $I$-adically complete
$R$-module and let $N \subset M$ be a submodule. The following are equivalent
\begin{enumerate}
\item $N$ is closed in $M$,
\item $N = \bigcap_{n \geq 1} (N + I^n M)$,
\item $M/N$ is $I$-adically complete
\end{enumerate}
If $I$ is finitely generated, these conditions imply that $N$
is $I$-adically complete.
\end{lemma}

\begin{proof}
The equivalence of (1) and (2) follows from
Formal Spaces, Lemma \ref{formal-spaces-lemma-closed}.
The equivalence of (2) and (3) is
Algebra, Lemma \ref{algebra-lemma-quotient-complete}.
Assume $I$ is finitely generated and (1) holds.
Let $N^\wedge$ be the $I$-adic completion of $N$. Since $M$
is $I$-adically complete we have a factorization
$$
N \to N^\wedge \to M
$$
of the inclusion map. Since $N^\wedge \to M$ is continuous and
$N$ is dense in $N^\wedge$ and $N$ is closed in $M$, we see that
$N^\wedge \to M$ factors through $N$. Thus $N^\wedge = N \oplus C$
for some $A$-module $C$. Thus $N$ is a direct summand of the
$I$-adically complete module $N^\wedge$, see
Algebra, Lemma \ref{algebra-lemma-hathat-finitely-generated}
(this is where we use that $I$ is finitely generated),
and hence $N$ is $I$-adically complete.
\end{proof}

\end{document}
